\begin{table}[htbp]
\centering\small
\begin{tabularx}{\textwidth}{Xlll}
\toprule
Quantity & A & B & C2\\
\midrule
Fresh customer reports accepted & 4/4 & 4/4 & 3/4\\
Recurring model/curator stage & Recorded & Recorded & None in execution\\
Complete recurring dollar cost & Unknown & Unknown & Unknown\\
Complete incremental construction & No new asset & Unknown & Unknown\\
Matched active customer/operator minutes & Unknown & Unknown & Unknown\\
Actual customer collections & \$0 & \$0 & \$0\\
\bottomrule
\end{tabularx}
\caption{The zero runtime model stage in C2 cannot earn an economic discount while complete equivalent work fails. Authoring a conventional program still incurs construction resources and unknown charges.}
\end{table}

The canonical shadow-ledger projection retains unknown total cost, profit and available operating surplus; it does not convert measured process resources into money. The final accounting snapshot gives one identity per recorded command, includes unsuccessful events, and preserves earlier snapshots and unknown discovery/authoring charges. The B and C2 asset manifests bind the actual tested entry points and source closures. Inactive comparison files included by provenance checks receive no productive credit.

For illustration only, a genuinely complete \$5 accepted job would break even against the developer wage anchor at approximately 4.59 active minutes, or against the QA wage anchor at 5.98 minutes. Neither the \$5 total nor those human durations was measured here. Likewise, a complete \$100 investment with a \$5 saving per equivalent accepted job and no upkeep would require twenty such jobs by simple arithmetic; those savings are not established for B or an accepted C2 replacement. The simulation must separately account for utilization, competition, failed work, maintenance, payment delays and foregone delivery.
